\section{Necesidades Detectadas}

Las necesidades detectadas sintetizan la fase de empatía del proceso de diseño: traducen lo observado en encuestas, entrevistas, observación directa y mapas de empatía en requerimientos que el rediseño debe atender. No describen solo funcionalidades ausentes, sino brechas entre la experiencia actual de WebSIS y la experiencia que el estudiante necesita para gestionar su trayectoria académica con autonomía y confianza.


{\scriptsize
\setlength{\tabcolsep}{3pt}
\renewcommand{\arraystretch}{1.1}

\begin{longtable}{|>{\bfseries}p{0.6cm}|p{2.6cm}|p{4.1cm}|p{3.7cm}|p{3.7cm}|}
\hline
\# & Necesidad & Descripción & Necesidad emocional subyacente & Evidencia \\
\hline
\endfirsthead

\hline
\# & Necesidad & Descripción & Necesidad emocional subyacente & Evidencia \\
\hline
\endhead

N1 & Guía en inscripción & El sistema debe guiar paso a paso el proceso de inscripción con indicadores de progreso visibles. & Seguridad y control: el estudiante necesita sentir acompañamiento durante la tarea. & 76\% encontró difícil la inscripción. E1 tardó 40 min con ayuda. S3 no completó la tarea. \\
\hline

N2 & Confirmación de acciones & El sistema debe confirmar claramente cada acción realizada para evitar incertidumbre. & Cierre y certeza: el usuario necesita saber que su acción fue registrada correctamente. & 56\% quedó con dudas tras completar tareas. E3 y E4 revisaban varias veces la pantalla. \\
\hline

N3 & Mensajes comprensibles & Los mensajes deben explicar qué ocurrió y cómo solucionarlo en lenguaje claro. & Confianza institucional: el usuario necesita entender los errores para actuar correctamente. & E4 perdió una materia por un mensaje confuso. 40\% reportó pasos no obvios. \\
\hline

N4 & Compatibilidad móvil & La interfaz debe funcionar correctamente en dispositivos móviles. & Equidad de acceso: muchos estudiantes dependen exclusivamente del celular. & E2: “Desde el celular es imposible”. 24\% usa celular como dispositivo principal. \\
\hline

N5 & Navegación alineada al modelo mental & El menú debe organizarse según las tareas del estudiante y no según la lógica interna del sistema. & Autonomía cognitiva: el usuario no debería aprender cómo funciona internamente el sistema. & S1 exploró 4 secciones antes de encontrar el módulo correcto. \\
\hline

N6 & Acceso rápido a tareas frecuentes & Las tareas más usadas deben estar accesibles desde la pantalla principal. & Eficiencia y ahorro de tiempo en períodos críticos. & 96\% usa inscripción y 92\% consulta notas frecuentemente. \\
\hline

N7 & Estabilidad técnica & El sistema debe mantenerse estable en períodos de alta demanda. & Confianza en la infraestructura del sistema. & E5 perdió una materia por fallas del sistema. \\
\hline

N8 & Soporte contextual integrado & El sistema debe incluir ayuda y orientación contextual dentro de cada módulo. & Autonomía institucional: reducir dependencia de terceros. & 80\% necesitó ayuda externa. 32\% pidió tutorial integrado. \\
\hline

N9 & Autonomía académica & El sistema debe permitir gestionar trámites sin depender constantemente de otras personas. & Dignidad y confianza en la propia capacidad de gestión. & 64\% siente que el sistema no brinda información suficiente para actuar solo. \\
\hline

\caption{Necesidades detectadas en usuarios de WebSIS}
\label{tab:necesidades-detectadas}
\end{longtable}
}
\subsection {Relación entre necesidades y perfiles de usuario}

Las necesidades detectadas no afectan de igual manera a los dos perfiles identificados. Algunas resultan especialmente críticas para el estudiante ingresante, mientras que otras impactan de forma transversal tanto a usuarios novatos como regulares. Esta diferenciación permite priorizar intervenciones de diseño según el nivel de vulnerabilidad y las consecuencias académicas asociadas a cada perfil.

{\small
\setlength{\tabcolsep}{4pt}
\renewcommand{\arraystretch}{1.2}
\begin{longtable}{|>{\bfseries}p{3.1cm}|p{5.35cm}|p{5.35cm}|}
\hline
Necesidad & Ingresante & Regular \\
\hline
\endfirsthead

\hline
Necesidad & Ingresante & Regular \\
\hline
\endhead

N1 -- Guía en inscripción & Crítica: no puede completar la tarea sin ayuda externa. La ansiedad anticipatoria es máxima porque cualquier error puede afectar su semestre. & Moderada: conoce parcialmente la ruta, pero debe reconstruirla cada semestre. \\
\hline

N2 -- Confirmación de acciones & Crítica: necesita certeza para cerrar emocionalmente el proceso y confiar en el resultado de sus acciones. & Crítica: revisa múltiples veces las acciones por experiencias previas de errores silenciosos. \\
\hline

N3 -- Mensajes comprensibles & Crítica: no tiene experiencia suficiente para interpretar mensajes ambiguos o técnicos. & Crítica: errores incomprensibles ya generaron consecuencias académicas reales. \\
\hline

N4 -- Compatibilidad móvil & Crítica: el celular suele ser su dispositivo principal de acceso. & Crítica: para usuarios sin laptop personal, el celular es el único canal disponible. \\
\hline

N5 -- Arquitectura de navegación & Crítica: navega por ensayo y error durante tareas fundamentales. & Moderada: ha memorizado rutas no documentadas, pero la estructura continúa siendo confusa. \\
\hline

N6 -- Acceso rápido a tareas frecuentes & Crítica: cada sesión comienza desde cero sin puntos de entrada claros. & Moderada: las rutas memorizadas compensan parcialmente la falta de accesos directos. \\
\hline

N7 -- Estabilidad técnica & Crítica: posee menor capacidad para interpretar y reaccionar ante fallos del sistema en momentos críticos. & Crítica: ha normalizado las caídas, pero continúa sufriendo consecuencias académicas reales. \\
\hline

N8 -- Soporte contextual integrado & Crítica: depende completamente de orientación externa durante sus primeras experiencias. & Moderada: lo requiere principalmente para tareas nuevas o poco frecuentes. \\
\hline

N9 -- Autonomía académica del estudiante & Crítica: la dependencia desde el primer contacto afecta su confianza y sensación de pertenencia institucional. & Crítica: incluso después de años de uso, muchos estudiantes aún sienten que el sistema no les proporciona suficiente información para actuar de forma autónoma. \\
\hline

\caption{Relación entre necesidades detectadas y perfiles de usuario}\label{tab:relacion-necesidades-perfiles}
\end{longtable}
}

